\documentclass{article}

\usepackage{listings}
\usepackage[english]{babel}


\usepackage[letterpaper,top=2cm,bottom=2cm,left=3cm,right=3cm,marginparwidth=1.75cm]{geometry}


\usepackage{amsmath}
\usepackage{graphicx}
\usepackage[colorlinks=true, allcolors=blue]{hyperref}

\title{Labwork 1 Report}
\author{Truong Quang Nam}


\begin{document}
\maketitle



\section{Introduction}



Gradient descent is a mathematical technique that iteratively finds the weights and bias that produce the model with the lowest loss. Gradient descent finds the best weight and bias by repeating the following process for a number of user-defined iterations

\section{Implementation}

\subsection{Gradient Descent Formula}

X = X - r x f'(X)

X: input
r: learning rate
f(X): output 
f'(X): derivative of f(x)

\subsection{Implement on $f(X) = X^2$}
In this step, we define the function to compute f(X) and it's derivative
\begin{lstlisting}
def y(x):
    return x*x
def y_p(x):
    return 2*x  

\end{lstlisting}


After that, I apply the formula to display update X make f(X) smaller by using gradient descent, update X until f(X) \textless  0.01

\begin{lstlisting}
def gradient_descent(x,lr):
    time = 0
    while y(x) > 0.01:
        x = x - lr * y_p(x)
        time = time + 1
        print(f"Epoch {time} x:{x:.3f} f(x):{y(x):.3f}")

\end{lstlisting}
\subsection{Result on different learning rate}
\begin{table}[h!]
\centering
\caption{Optimization Results (Learning Rate = 0.1)}
\label{tab:optimization_results}
\begin{tabular}{|c|c|c|}
\hline
\textbf{Epoch} & \textbf{x} & \textbf{f(x)} \\ \hline
1  & 8.000 & 64.000 \\ \hline
2  & 6.400 & 40.960 \\ \hline
3  & 5.120 & 26.214 \\ \hline
4  & 4.096 & 16.777 \\ \hline
5  & 3.277 & 10.737 \\ \hline
6  & 2.621 & 6.872  \\ \hline
7  & 2.097 & 4.398  \\ \hline
8  & 1.678 & 2.815  \\ \hline
9  & 1.342 & 1.801  \\ \hline
10 & 1.074 & 1.153  \\ \hline
11 & 0.859 & 0.738  \\ \hline
12 & 0.687 & 0.472  \\ \hline
13 & 0.550 & 0.302  \\ \hline
14 & 0.440 & 0.193  \\ \hline
15 & 0.352 & 0.124  \\ \hline
16 & 0.281 & 0.079  \\ \hline
17 & 0.225 & 0.051  \\ \hline
18 & 0.180 & 0.032  \\ \hline
19 & 0.144 & 0.021  \\ \hline
20 & 0.115 & 0.013  \\ \hline
21 & 0.092 & 0.009  \\ \hline
\end{tabular}
\end{table}
\begin{table}[h!]
\centering
\caption{Optimization Results (Learning Rate = 0.2)}
\label{tab:optimization_02}
\begin{tabular}{|c|c|c|}
\hline
\textbf{Epoch} & \textbf{x} & \textbf{f(x)} \\ \hline
1  & 6.000 & 36.000 \\ \hline
2  & 3.600 & 12.960 \\ \hline
3  & 2.160 & 4.666  \\ \hline
4  & 1.296 & 1.680  \\ \hline
5  & 0.778 & 0.605  \\ \hline
6  & 0.467 & 0.218  \\ \hline
7  & 0.280 & 0.078  \\ \hline
8  & 0.168 & 0.028  \\ \hline
9  & 0.101 & 0.010  \\ \hline
10 & 0.060 & 0.004  \\ \hline
\end{tabular}
\end{table}

\end{document}
